\chapter{De Analista a Diseñador: La Simulación como Herramienta de Optimización}
\label{chap:optimization}

%==================================================================================
\section{El Siguiente Nivel: De la Predicción al Diseño}
%==================================================================================

A lo largo de este libro, hemos construido un poderoso arsenal de herramientas para predecir el comportamiento de sistemas físicos complejos. Hemos respondido a la pregunta: ``Dado este diseño, ¿cómo se comportará?''. Sin embargo, la pregunta fundamental en ingeniería es a menudo la inversa: ``Dado este comportamiento deseado, ¿cuál es el mejor diseño posible?''.

Este capítulo final cierra el círculo de nuestro viaje. Pasamos del rol de analista al de diseñador. Introducimos el campo de la \textbf{Optimización Basada en Simulación}, donde unimos nuestros solvers de elementos finitos con algoritmos de optimización matemática para crear un bucle que no solo predice, sino que mejora sistemáticamente un diseño.

En lugar de realizar un proceso manual de prueba y error, automatizaremos la búsqueda del diseño óptimo. Esto representa el pináculo del uso de la simulación en la industria, permitiendo la creación de componentes más ligeros, más eficientes y más robustos de lo que sería posible mediante la intuición humana solamente.

\begin{center}
\begin{tabular}{|p{0.9\textwidth}|}
\hline
\textbf{CONTEXTO INGENIERIL} \\
La optimización basada en simulación está revolucionando el diseño industrial:
\begin{itemize}
    \item \textbf{Industria Aeroespacial y Automotriz:} El \textbf{diseño generativo} y la \textbf{optimización topológica} se utilizan para crear componentes estructurales (soportes, chasis) que son extremadamente ligeros y rígidos, con formas orgánicas que a menudo solo pueden ser fabricadas mediante impresión 3D.
    \item \textbf{Diseño de Disipadores de Calor:} Optimizar la forma y disposición de las aletas de un disipador para maximizar la transferencia de calor con la mínima cantidad de material.
    \item \textbf{Diseño de Implantes Médicos:} Crear implantes (ej. de cadera) con una distribución de material que minimice los esfuerzos en el hueso circundante y maximice la vida útil del implante.
    \item \textbf{Ingeniería de Procesos:} Optimizar los parámetros operativos de un reactor químico para maximizar la producción de un compuesto deseado.
\end{itemize}
\\
\hline
\end{tabular}
\end{center}

%==================================================================================
\section{La Anatomía de un Problema de Optimización}
\label{sec:opt-theory}
%==================================================================================
Cualquier problema de optimización, sin importar su complejidad, puede ser descompuesto en tres componentes fundamentales.

%----------------------------------------------------------------------------------
\subsection{La Función Objetivo (\(J\))}
%----------------------------------------------------------------------------------
La función objetivo es la \textbf{métrica cuantificable del rendimiento} que queremos mejorar. Es una única cantidad escalar que nuestro algoritmo intentará minimizar o maximizar. La elección de una buena función objetivo es el paso más crítico en la formulación de un problema de optimización.
\begin{itemize}
    \item \textbf{Mecánica Estructural:} Minimizar la masa del componente; minimizar el desplazamiento máximo; minimizar la concentración de esfuerzos.
    \item \textbf{Transferencia de Calor:} Maximizar la tasa de transferencia de calor; minimizar la temperatura máxima.
    \item \textbf{Dinámica de Fluidos:} Minimizar la caída de presión (resistencia aerodinámica); maximizar el coeficiente de sustentación.
\end{itemize}
La función objetivo es a menudo una integral sobre el dominio, por ejemplo, la energía de deformación elástica del cuerpo: \(J = \int_{\Omega} \bm{\sigma} : \bm{\varepsilon} \, dV\).

%----------------------------------------------------------------------------------
\subsection{Las Variables de Diseño (\(\alpha\))}
%----------------------------------------------------------------------------------
Las variables de diseño son los \textbf{parámetros que podemos cambiar} para intentar mejorar la función objetivo. Representan la ``libertad'' que tiene el optimizador para modificar el diseño.
\begin{itemize}
    \item \textbf{Optimización de Forma:} Las coordenadas de los puntos de control de una spline que define el contorno de un objeto.
    \item \textbf{Optimización Paramétrica:} El espesor de una viga, el radio de un agujero, la temperatura de un proceso.
    \item \textbf{Optimización Topológica:} La más poderosa de todas. La variable de diseño es la presencia o ausencia de material en cada punto del dominio.
\end{itemize}

%----------------------------------------------------------------------------------
\subsection{Las Restricciones (\(g\))}
%----------------------------------------------------------------------------------
Las restricciones son los \textbf{límites y reglas del juego} que el diseño debe respetar. Sin restricciones, la mayoría de los problemas de optimización tienen soluciones triviales (ej. para minimizar la masa, la solución es no tener material).
\begin{itemize}
    \item \textbf{Restricción de Volumen/Masa:} El volumen total de material no puede exceder un cierto porcentaje del volumen de diseño disponible.
    \item \textbf{Restricciones de Esfuerzo/Deformación:} El esfuerzo máximo en cualquier punto del componente no puede exceder el límite elástico del material. El desplazamiento máximo no puede superar un umbral.
    \item \textbf{Restricciones Geométricas:} Ciertas superficies deben permanecer fijas; el espesor mínimo de una pared no puede ser menor a un valor fabricable.
\end{itemize}

El problema de optimización formal se escribe entonces como:
\begin{equation}
    \min_{\alpha} J(\bm{u}(\alpha)) \quad \text{sujeto a} \quad g_i(\bm{u}(\alpha)) \le 0
\end{equation}
donde \(\bm{u}(\alpha)\) es la solución de la EDP (nuestro solver FEM) para un conjunto dado de variables de diseño \(\alpha\).

%==================================================================================
\section{Optimización Topológica: El Arte de Encontrar la Forma}
\label{sec:topo-opt}
%==================================================================================
La \textbf{optimización topológica} es una de las áreas más espectaculares de la optimización basada en simulación. En lugar de simplemente modificar una forma preexistente, comienza con un bloque de material y ``esculpe'' la estructura óptima eliminando el material que no es estructuralmente eficiente.

El método más popular es el \textbf{Método SIMP (Solid Isotropic Material with Penalization)}, detallado en la obra seminal de Bendsøe y Sigmund \cite{Bendsoe2003}. La idea es asignar una variable de diseño de densidad, \(\rho_e \in [0,1]\), a cada elemento de la malla.
\begin{itemize}
    \item Si \(\rho_e = 1\), el elemento es de material sólido.
    \item Si \(\rho_e = 0\), el elemento es un vacío.
    \item Valores intermedios son penalizados para empujar la solución final a un diseño ``blanco y negro''.
\end{itemize}
El optimizador comienza con un bloque de material (\(\rho_e=1\) en todas partes) y, en cada iteración, resuelve el problema de elasticidad, calcula la sensibilidad de la función objetivo a la densidad de cada elemento, y utiliza esta información para redistribuir el material, quitándolo de las zonas de bajo esfuerzo y poniéndolo en las de alto esfuerzo, todo ello mientras respeta una restricción de volumen global.

El resultado son diseños de una eficiencia asombrosa, a menudo con formas orgánicas y complejas que un diseñador humano difícilmente podría concebir. La optimización topológica es el verdadero motor detrás del \textbf{diseño generativo} y representa la sinergia perfecta entre la simulación, la optimización y las nuevas tecnologías de fabricación como la impresión 3D.